% Edition 0.5. Source snapshot d4870eb1, exact original normalized Fock observables.
\chapter[The Normalized Fock-Field Program]{The Normalized Original Fock-Field Program: Bounded Observables and Genuine $L^2$ Classes}
\label{ch:norm}

\section{Why we began another operator program after constructing the Hamiltonian}
The original physical Hamiltonian of Chapter~\ref{ch:hamiltonian} is an
unbounded operator on a gauge-invariant Hilbert space. It is tempting to
read any subsequent bounded-operator theorem as a bound for the original
fermionic interaction. That interpretation would be false. Our NORM-0
program addresses a separate, reusable question: can a particular
configuration-dependent quadratic expression in the \emph{original}
Majorana operators be normalized into a uniformly bounded field, and can
its pointwise action be constructed on the actual vector-valued Lebesgue
$L^2$ space without changing representatives or carriers?

The historical reason for the small-gate design is instructive. A bounded
finite-dimensional operator at every $X$ need not be a valid multiplier
unless measurability and almost-everywhere equivalence are handled; nor
may a fixed-operator \texttt{compLp} theorem be applied to an
$X$-dependent operator. Our source-first construction proceeds in
four distinct steps: an algebraic Gram bound, a continuous field,
$L^2$ membership of pointwise images, and quotient-level representative
independence. Each is a different proof obligation.

\section[The actual spatial Gram matrix and the fixed-B field]{The actual spatial Gram matrix and the fixed-$B$ field}
For $X\in\mathbb R^{9\times3}$, let $x_i\in\mathbb R^3$ denote
its $i$th color row, and write
\begin{equation}\label{eq:norm-gram}
  M_{ij}(X)=\langle x_i,x_j\rangle_{\mathbb R^3},
  \qquad \tau(X)=\sum_{i=0}^{8}M_{ii}(X)=\|X\|^2.
\end{equation}
The equality with the squared Euclidean norm is notation for the
\emph{original} source spatial Gram trace, not an identification with the
quartic potential. The Cauchy--Schwarz and positivity of a Gram matrix
supply
\begin{equation}\label{eq:norm-gram-entry-bound}
   |M_{ij}(X)|\le\sqrt{M_{ii}(X)M_{jj}(X)}\le\tau(X).
\end{equation}
The source qualified NORM-0A matrix bound is the exact finite-index
version of this elementary argument.

Let $\mathcal F$ denote the literal original \texttt{Fermion 2}, and let
$B\in\operatorname{Mat}_{16}(\mathbb R)$ be a \emph{fixed} real matrix.
For the original nine gamma matrices $\gamma_i$ and the already
constructed color-contracted quadratic Fock-operator map
$\mathcal Q:\operatorname{Mat}_{16}(\mathbb R)\to\operatorname{End}_{\mathbb C}(\mathcal F)$,
put
\begin{equation}\label{eq:norm-unscaled-fock}
 A_B(X)=\sum_{i,j=0}^{8}M_{ij}(X)\,\mathcal Q(\gamma_iB\gamma_j),
 \qquad
 C_B=\sum_{i,j=0}^{8}\|\mathcal Q(\gamma_iB\gamma_j)\|_{\rm op}.
\end{equation}
The source writes $A_B$ as \texttt{originalFockGramObservable};
$\mathcal Q$ is the \emph{original} color-quadratic operator map,
not a newly chosen abstract Clifford representation.

\begin{theorem}[Uniformly bounded normalized original Fock field;
strict-qualified]\label{thm:norm-field}
Define, on the genuine original complex Fock carrier,
\begin{equation}\label{eq:norm-field}
 \widetilde A_B(X)=\frac{i}{1+\tau(X)}A_B(X).
\end{equation}
Then $C_B\ge0$, $1+\tau(X)>0$, and
\begin{equation}\label{eq:norm-field-bound}
  \|A_B(X)\|_{\rm op}\le\tau(X)C_B,
 \qquad \|\widetilde A_B(X)\|_{\rm op}\le C_B
 \quad\text{for all }X.
\end{equation}
This is a statement about bounded Fock \emph{fiber operators} and a
fixed $B$; it does not assert self-adjointness for arbitrary $B$.
\end{theorem}
\begin{proof}[The source finite-sum proof]
Apply the triangle inequality to the $9\times9$ sum in
\eqref{eq:norm-unscaled-fock}, the operator norm identity for scalar
multiplication, and \eqref{eq:norm-gram-entry-bound}:
\[
  \|A_B(X)\|\le\sum_{ij}|M_{ij}(X)|
  \|\mathcal Q(\gamma_iB\gamma_j)\|
  \le\tau(X)C_B.
\]
Since $\tau(X)\ge0$, division by $1+\tau(X)$ gives the second
inequality. The factor $i$ has modulus one. Every operator here acts
on the already fixed original $2^{24}$-dimensional Fock space; neither
the proof nor the normalization changes the physical Hamiltonian.
\end{proof}
\formalrecord{\texttt{BFSSNORM0COriginalNormalizedFockField.lean},
\texttt{originalFockGammaBoundConstant\_nonneg},
\texttt{originalFockGram\_norm\_le\_trace},
\texttt{originalNormalizedFockGram\_norm\_le}; source blob
\texttt{84de66e73be2018b03ced898ebf70e691c07b40f};
\href{https://github.com/etblink/math/actions/runs/43017992525}{strict run \#43017992525}.
Atlas N-01.}

\section{Continuity, measurable evaluation, and the non-finite measure trap}
The underlying bosonic space is $\mathbb R^{27}$ with ordinary
Lebesgue volume. It is not a probability space, and its total volume
is infinite. In particular, the constant function with values in a
nonzero finite-dimensional Fock operator space does not become
$L^2$ merely because its values are uniformly bounded.

The valid conclusion instead concerns \emph{products with an $L^2$
argument}. For fixed $B$, the source continuity gate establishes
$X\mapsto\widetilde A_B(X)$ as a continuous function into the
bounded-operator norm topology. For $f\in\OriginalLtwo=L^2(\mathbb R^{27};\mathcal F)$
with representative $f(X)$, the evaluation map
$(S,v)\mapsto Sv$ is continuous bilinear. It follows that
$X\mapsto\widetilde A_B(X)f(X)$ is strongly measurable almost
everywhere. The domination inequality
\begin{equation}\label{eq:norm-lp-domination}
  \|\widetilde A_B(X)f(X)\|\le C_B\|f(X)\|
\end{equation}
then gives the $L^2$ integrability of the output without requiring
finite total measure.

\begin{theorem}[Pointwise original $L^2$ membership;
strict-qualified]\label{thm:norm-memlp}
For every original $f\in\texttt{FullL2 2}$ and fixed real matrix $B$,
$X\mapsto\widetilde A_B(X)f(X)$ is a.e. strongly measurable and is
\texttt{MemLp} of exponent $2$ on the actual source bosonic Lebesgue
measure. Moreover, \eqref{eq:norm-lp-domination} holds pointwise for
the source representative.
\end{theorem}
\begin{proof}[Accepted bounded-evaluation and domination argument]
The operator field's source continuity gives its a.e. strong
measurability. The built-in source $L^2$ class supplies the a.e.
strong measurability of $f$. Applying the pinned Mathlib continuous
bilinear-evaluation interface yields measurable output; the NORM-0C
norm estimate supplies the domination. \texttt{MemLp.of\_le\_mul}
then proves membership, and its hypotheses do not require a finite
measure of the full bosonic domain.
\end{proof}
\formalrecord{\texttt{BFSSNORM0D1OriginalFockFieldContinuity.lean} and
\path{BFSSNORM0D2OriginalPointwiseL2Membership.lean}, including
\path{originalPointwiseFock_aeStronglyMeasurable},
\texttt{originalPointwiseFock\_norm\_le}, and
\texttt{originalPointwiseFock\_memLp}; strict D1-A receipt
\href{https://github.com/etblink/math/actions/runs/43047181537}{\#43047181537}.
Atlas N-02.}

\section{From a pointwise function to an actual almost-everywhere class}
An $L^2$ element is an equivalence class modulo null sets, not a
canonical pointwise wavefunction. One must be careful when saying a
multiplication rule is independent of representatives.

\begin{lemma}[Representative-safe $L^2$ construction;
strict-qualified]\label{lem:norm-carrier}
Let $q:\mathbb R^{27}\to\mathcal F$ be a function with an actual
\texttt{MemLp} certificate. The original \texttt{MemLp.toLp}
construction yields an element $[q]\in\OriginalLtwo$ whose
representative agrees with $q$ almost everywhere. If $q=q'$ almost
everywhere and both have the required integrability, their $L^2$
classes coincide.
\end{lemma}
\begin{proof}
This is the exact quotient property of the source carrier used by the
NORM-0D3A bridge. It uses \texttt{MemLp.coeFn\_toLp} and
\texttt{MemLp.toLp\_congr}; changing the function on a null set
changes neither its class nor its norm.
\end{proof}
\formalrecord{\texttt{BFSSNORM0D3AOriginalLpCarrierBridge.lean},
\texttt{originalCarrierLpOfMemLp},
\texttt{originalCarrierLpOfMemLp\_ae}, and
\texttt{originalCarrierLpOfMemLp\_congr}; blob
\texttt{3fc16b042fe51f1e1a8c52260d5bba5135831ffc};
\href{https://github.com/etblink/math/actions/runs/43084968990}{strict run \#43084968990}.
Atlas N-03.}

The next gate applies this carrier construction to the actual
normalized Fock image and proves the almost-everywhere output law.
The accepted theorem deliberately avoids a new named multiplication
wrapper, because an earlier attempt at a large instantiated wrapper
encountered Lean recursion. That failure does not falsify the intended
mathematical operator; it changes how the result is packaged and
verified.
\formalrecord{\texttt{BFSSNORM0D3CDirectOriginalFockLpOutput.lean},
\texttt{originalFockLpDirect\_ae}; blob
\texttt{bd9146fe825b4675bfb62f772bc02f83e204d30b};
\href{https://github.com/etblink/math/actions/runs/43091441892}{strict run \#43091441892}.
Atlas N-04.}

\section{A short mathematical deduction and an important qualification hold}
The accepted results suggest the natural formula
\begin{equation}\label{eq:norm-multiplier-definition}
   (\mathsf M_B f)(X):=\widetilde A_B(X)f(X)
   \qquad\text{a.e.},
\end{equation}
which, in ordinary functional analysis, describes a bounded linear
multiplication operator on $L^2$ satisfying
\begin{equation}\label{eq:norm-multiplier-goal}
   \|\mathsf M_B f\|_2\le C_B\|f\|_2.
\end{equation}
\begin{proposition}[Bounded multiplier from the source normalized
field; conventional consequence]\label{prop:normalized-class-multiplier}
For each \emph{fixed} real matrix $B$, the field
$\widetilde A_B$ determines a bounded complex-linear operator
$\mathsf M_B:\OriginalLtwo\to\OriginalLtwo$ in ordinary
Hilbert-space mathematics, with
\begin{equation}\label{eq:norm-multiplier-consequence}
 \|\mathsf M_B\|_{\mathrm{op}}\le C_B.
\end{equation}
Its action on any representative is
$\widetilde A_B(X)f(X)$ almost everywhere.
\end{proposition}
\begin{proof}
The source continuity result makes $\widetilde A_B$ strongly
measurable, and the exact NORM-0C theorem supplies the uniform
operator bound $\|\widetilde A_B(X)\|\le C_B$.
Proposition~\ref{prop:bounded-l2-multiplier}
therefore applies to the original Fock-vector $L^2$ carrier:
pointwise evaluation preserves almost-everywhere equality,
the output is square integrable, and multiplication
is complex-linear with the stated norm bound.
No finite-volume assumption is used.
\end{proof}
This is a \emph{conventional} mathematical deduction from the
displayed source inputs, not a newly checked source Lean
\texttt{ContinuousLinearMap} theorem. In particular, it
does not promote the next representative-transport source
\emph{gate} to strict-qualified status or identify
$\mathsf M_B$ with the original unbounded fermionic interaction.

The next exact source candidate makes the representative-independence
step explicit: if $g=f$ almost everywhere, then
$\widetilde A_B(X)g(X)=\widetilde A_B(X)f(X)$ almost everywhere, and
the former is $\texttt{MemLp}$ whenever the latter is. Its source
proof is small, but the required independently completed strict gate
\href{https://github.com/etblink/math/actions/runs/43095354808}{\#43095354808}
was still \emph{queued with no conclusion} at our 11 November
editorial source check. It must not be counted as a green qualification
without a complete run and consumer audit.
\formalrecord{Pending source
\texttt{BFSSNORM0D3DOriginalFockRepresentativeTransport.lean},
\texttt{originalFockImage\_ae\_of\_rep} and
\texttt{originalFockImage\_memLp\_of\_rep}; candidate blob
\texttt{cb16f081fb1dbff0f149e94b39145d3ea2dec58b}.
Atlas N-05, \textbf{pending strict acceptance}.}

\section{The critical scientific boundary}
The observable $\widetilde A_B$ is normalized by $1+\|X\|^2$;
its unnormalized numerator is quadratic in $X$. The physical
fermionic interaction $B_M^{(1,0)}(X)$ of
Chapter~\ref{ch:energy}, however, has a different algebraic
coefficient construction and grows \emph{linearly} in $X$. A uniform
bound on $\widetilde A_B$ proves neither the desired source linear
growth inequality \eqref{eq:B-linear-target} nor physical-sector
coercivity. Any comparison must be a separately typed source theorem
identifying the actual operators, with the signs, indices, domains and
normalization intact. Without it, the two lanes are useful but distinct.

The positive role of NORM-0 is methodological and mathematical: it
exhibits how to build physically sourced, gauge-compatible bounded
operator observables on an infinite-volume $L^2$ carrier, with rigorous
handling of null-set representatives. This infrastructure may aid
future compactness estimates or gauge-equivariant operator calculus;
its usefulness is not evidence of an eigenpair or a mass gap.
